\section{Theoretical Foundations}\label{sec:theory}
This section reviews the mathematical foundations of classical and quantum clustering, including the definition of the 5D color-spatial manifold, parallel Single Instruction, Multiple Threads (SIMT) GPU execution complexity, and the mathematical derivation of the quantum SWAP test phase mapping.

\subsection{Classical K-Means and the 5D Color-Spatial Manifold}
Classical K-Means groups $N$ data points into $K$ uniform clusters \cite{lloyd1982least} by iteratively minimizing WCSS \cite{krzanowski1988criterion}:
\begin{equation}
\label{eq:wcss}
J = \sum_{i=1}^{K} \sum_{\mathbf{x} \in S_i} \|\mathbf{x} - \boldsymbol{\mu}_i\|^2
\end{equation}
where $S_i$ is the set of points in the $i$-th cluster, and $\boldsymbol{\mu}_i$ is the centroid of $S_i$. Each iteration consists of an assignment step based on Euclidean distance and an update step that calculates the average position of points assigned to each cluster \cite{lloyd1982least}.

Traditional clustering operates solely in 3D color spaces (e.g., RGB or CIELAB) \cite{gonzalez2009digital}, which can lead to spatial fragmentation. Spatial consistency is enforced by extending the feature space to a 5D manifold \cite{achanta2012slic}:
\begin{equation}
\mathbf{x} = [B, G, R, x, y]^T
\end{equation}
where $B, G, R$ represent color channels, and $x, y$ are pixel coordinates. A spatial weight parameter ($\lambda_{\text{spatial}}$) is introduced into the distance metric to balance the different scales of color values and spatial coordinates \cite{shi2000normalized}:
\begin{align}
\|\mathbf{x}_n - \boldsymbol{\mu}_i\|^2 &= (B_n - \boldsymbol{\mu}_{i,B})^2 + (G_n - \boldsymbol{\mu}_{i,G})^2 \nonumber \\
&{}+ (R_n - \boldsymbol{\mu}_{i,R})^2 + \lambda_{\text{spatial}} \bigl[ (x_n - \boldsymbol{\mu}_{i,x})^2 \nonumber \\
&{}+ (y_n - \boldsymbol{\mu}_{i,y})^2 \bigr]
\end{align}
In a continuous video stream, the centroids change gradually between frames. This temporal redundancy is exploited via a warm-start strategy \cite{szeliski2022computer}, where the converged centroids from the previous frame serve as the starting seeds for the current frame, reducing convergence iterations.

Clustering quality is evaluated using three standard metrics: WCSS, the Davies-Bouldin ($\mathcal{D}\mathcal{B}$) Index \cite{davies1979cluster}, and the Silhouette Coefficient \cite{rousseeuw1987silhouettes}. The $\mathcal{D}\mathcal{B}$ index measures the ratio of cluster dispersion to centroid separation. The Silhouette Coefficient evaluates how well each point fits its cluster, but exactly calculating it for a VGA frame ($N=307,200$) requires $\approx 47.19$ billion distance calculations per frame, which is impossible in real-time. A Monte Carlo approximation bypasses this bottleneck by sampling a random, uniform subset of $M=2,000$ pixels, reducing complexity to $\mathcal{O}(M^2)$ \cite{shutaywi2021silhouette}.

\subsection{Parallel GPU Architectures and CUDA SIMT}
Running K-Means sequentially on a CPU is too slow for real-time video \cite{farivar2008parallel}. For a VGA frame with $K=8$ and $I=15$ iterations, the CPU must perform over $36.8$ million calculations per frame. Achieving $120\text{ FPS}$ \cite{cruden2019simulator, jerald2009scene} requires completing this computation in under $8.33\text{ ms}$, which is hindered by single-core bottlenecks \cite{cook2012cuda}.

Using NVIDIA's CUDA, calculations are offloaded to the GPU using the SIMT model \cite{CudaGuide}. Threads execute in groups of 32 (Warps), sharing low-latency registers and Streaming Multiprocessor (SM) caches \cite{cook2012cuda}. Maximizing execution throughput requires avoiding warp divergence (conditional branching) and global memory lookup stalls by utilizing coalesced memory transactions and fast on-chip shared memory (\texttt{\_\_shared\_\_}) \cite{CudaGuide, farivar2008parallel}.

CPU-bound preprocessing bottlenecks are eliminated by implementing a parallel strided downsampling preprocessor directly in the CUDA kernel. By scaling the GPU execution grid to the exact downsampled size $W_{\text{out}} = \lceil W/S \rceil$ and mapping threads directly to input array indices via the stride $S$, warp divergence is eliminated and execution latency remains deterministic.

\subsection{Quantum-Inspired Emulation and SWAP Test Circuit}
\subsubsection{QML and NISQ Physical Constraints}
While quantum machine learning (QML) offers theoretical speedups, current quantum co-processors exist in the NISQ era, characterized by high gate error rates and short coherence times ($T_1, T_2$) \cite{bharti2022noisy, preskill2018quantum}. Deep quantum circuits accumulate noise, rendering physical execution unreliable. Furthermore, live video streams suffer from remote QPU network latency \cite{fraunhofer2025quantum}, while QML algorithms themselves suffer from significant state preparation overhead when encoding classical features into qubits \cite{kerenidis2019q}. These limits are bypassed by simulating the quantum SWAP test deterministically using parallel CUDA kernels, mapping the mathematical state-overlap formula directly to classical GPU threads.

\subsubsection{SWAP Test Circuit and Mathematical Derivation}
In quantum information theory, the SWAP test measures the overlap (similarity) between two states $\lvert\psi\rangle$ and $\lvert\phi\rangle$ using three registers: an ancilla qubit initialized to $\lvert 0 \rangle_{\text{anc}}$, and two target registers holding the inputs \cite{buhrman2001quantum}. A Hadamard gate is applied to the ancilla, followed by a Controlled-SWAP (CSWAP) gate that swaps the targets if the ancilla is in state $\lvert1\rangle$, and a second Hadamard gate to induce interference.

The registers evolve as follows:
\begin{align}
\lvert\Psi_0\rangle &= \lvert 0 \rangle_{\text{anc}}\lvert\psi\rangle\lvert\phi\rangle \\
\lvert\Psi_1\rangle &= \frac{1}{\sqrt{2}}\lvert0\rangle_{\text{anc}}\lvert\psi\rangle\lvert\phi\rangle + \frac{1}{\sqrt{2}}\lvert1\rangle_{\text{anc}}\lvert\psi\rangle\lvert\phi\rangle \\
\lvert\Psi_2\rangle &= \frac{1}{\sqrt{2}}\lvert0\rangle_{\text{anc}}\lvert\psi\rangle\lvert\phi\rangle + \frac{1}{\sqrt{2}}\lvert1\rangle_{\text{anc}}\lvert\phi\rangle\lvert\psi\rangle
\end{align}
The final Hadamard gate results in:
\begin{equation}
\lvert\Psi_3\rangle = \frac{1}{2}\lvert0\rangle_{\text{anc}}(\lvert\psi\rangle\lvert\phi\rangle + \lvert\phi\rangle\lvert\psi\rangle) + \frac{1}{2}\lvert1\rangle_{\text{anc}}(\lvert\psi\rangle\lvert\phi\rangle - \lvert\phi\rangle\lvert\psi\rangle)
\end{equation}
The probability $P(\lvert 0 \rangle)$ of measuring the ancilla in the ground state is:
\begin{align}
P(\lvert0\rangle) &= \left\| \langle 0 \rvert \Psi_3 \rangle \right\|^2 = \left\| \frac{1}{2}(\lvert\psi\rangle\lvert\phi\rangle + \lvert\phi\rangle\lvert\psi\rangle) \right\|^2 \\
&= \frac{1}{4} \left[ \langle\psi\vert\psi\rangle\langle\phi\vert\phi\rangle + 2\lvert\langle\psi\vert\phi\rangle\rvert^2 + \langle\phi\vert\phi\rangle\langle\psi\vert\psi\rangle \right] \nonumber \\
&= \frac{1}{2} + \frac{1}{2}\lvert\langle\psi\vert\phi\rangle\rvert^2
\end{align}
This proves that the ground state probability is directly related to the state overlap (or fidelity) $\lvert\langle\psi\vert\phi\rangle\rvert^2$ \cite{buhrman2001quantum}.

\subsubsection{Hilbert Phase Space Mapping and Boundary Justification}\label{subsect:phasespace}\label{subsect:boundary_justification}
This state overlap is emulated directly in CUDA C++. Normalized 5D pixel features are mapped to phase states: $\theta_d = (p_d - c_d) \cdot \lambda_{\text{scale}}$, where $p_d$ is the pixel feature, $c_d$ is the centroid feature, and $\lambda_{\text{scale}}$ is a scaling factor. The 1D overlap is $\text{Overlap}_d = \cos^2(\theta_d)$. The global similarity is the product of these overlaps, representing the state overlap of the multi-qubit product state \cite{buhrman2001quantum}:
\begin{equation}
\text{Fidelity}(\mathbf{p}, \mathbf{c}) = \prod_{d=1}^{D} \cos^2\left( (p_d - c_d) \cdot \lambda_{\text{scale}} \right)
\end{equation}
Subtracting this from unity yields the quantum minimizing distance metric:
\begin{equation}
D_{\text{Quantum}}(\mathbf{p},\mathbf{c}) = 1.0 - \prod_{d=1}^{D} \cos^2\left( (p_d - c_d) \cdot \lambda_{\text{scale}} \right)
\end{equation}
The scaling coefficient is locked to $\lambda_{\text{scale}} = \frac{\pi}{4}$, bounding the angle within $[0, \frac{\pi}{4}]$, allowing the GPU to calculate the cosine quickly using hardware-accelerated instructions (\texttt{\_\_cosf}) \cite{CudaGuide}.

Unlike the classical Euclidean distance which scales linearly, the quantum-inspired metric changes along a squared cosine curve, resembling kernel-based projection methods that transform data into high-dimensional Hilbert spaces to resolve complex boundaries \cite{scholkopf1998nonlinear}. As the distance increases, the slope of $\cos^2(\theta)$ becomes steeper, yielding a sharp transition profile near the cluster borders that theoretically suppresses color-space bleeding and preserves edge transitions in noisy regions \cite{shi2000normalized}.
